% =========================================================
% TOPIC 1 -- INTRODUCTION
% =========================================================

\section{Introduction}

% =========================================================
% TOPIC TITLE SLIDE
% =========================================================

\begin{frame}{Introduction}

\begin{center}

\vspace{1cm}

{\Huge\textbf{Introduction}}

\vspace{0.6cm}

\textit{Applying fundamental concepts of mathematical reasoning to
computing scenarios}

\end{center}

\end{frame}

% =========================================================
% QUESTION 1 -- PART 1
% =========================================================

\begin{frame}{Introduction -- Question 1}

\begin{block}{Software Requirement Verification}

A software system has the following requirements:

\begin{itemize}
    \item Every registered user must have a unique user ID.
    \item Every user ID is associated with exactly one account.
    \item An inactive account must still have a unique user ID.
\end{itemize}

A developer claims:

\begin{center}
\textit{``If two accounts have the same user ID, they must belong to the same user.''}
\end{center}

\end{block}

\end{frame}

% =========================================================
% QUESTION 1 -- PART 2
% =========================================================

\begin{frame}{Introduction -- Question 1 (Continued)}

\begin{block}{Question}

Determine whether the developer's statement logically follows from the
given requirements.

\medskip

Explain your reasoning and identify the mathematical idea involved.

\medskip

Let

\[
P=\text{``two accounts have the same user ID.''}
\]

and let

\[
Q=\text{``the accounts belong to the same user.''}
\]

Express the developer's statement in logical form.

\end{block}

\end{frame}

% =========================================================
% ANSWER 1 -- PART 1
% =========================================================

\begin{frame}{Introduction -- Question 1: Answer}

\begin{exampleblock}{Solution}

The requirements state that every user ID is unique.

Therefore, two different accounts cannot have the same user ID.

Hence,

\[
P=\mathrm{False}.
\]

The developer's statement has the form

\[
P\rightarrow Q.
\]

\end{exampleblock}

\end{frame}

% =========================================================
% ANSWER 1 -- PART 2
% =========================================================

\begin{frame}{Introduction -- Question 1: Answer (Continued)}

\begin{exampleblock}{Solution}

Since the antecedent $P$ is false under the given requirements,

\[
P\rightarrow Q
\]

is logically true.

This illustrates the concept of \textbf{vacuous truth} in logical
implication.

\medskip

\textbf{Conclusion:}

The implication is formally true because its antecedent describes an
impossible situation. However, the requirements do not independently
establish the claimed relationship between the accounts.

\end{exampleblock}

\end{frame}

% =========================================================
% QUESTION 2 -- PART 1
% =========================================================

\begin{frame}{Introduction -- Question 2}

\begin{block}{AI Classification Rule}

An AI-based spam filter uses the following rule:

\begin{center}
\textit{A message is spam if it contains a suspicious link or more than
five promotional keywords.}
\end{center}

Consider a message containing a suspicious link but only two promotional
keywords.

Let

\[
P=\text{``The message contains a suspicious link.''}
\]

\end{block}

\end{frame}

% =========================================================
% QUESTION 2 -- PART 2
% =========================================================

\begin{frame}{Introduction -- Question 2 (Continued)}

\begin{block}{Question}

Let

\[
Q=\text{``The message contains more than five promotional keywords.''}
\]

\medskip

Determine whether the message is classified as spam.

\medskip

Express the spam rule using $P$ and $Q$.

\medskip

Why is knowing only the truth value of $Q$ insufficient to determine
whether the message is spam?

Explain your answer using the logical operator involved.

\end{block}

\end{frame}

% =========================================================
% ANSWER 2 -- PART 1
% =========================================================

\begin{frame}{Introduction -- Question 2: Answer}

\begin{exampleblock}{Solution}

The message contains a suspicious link.

Therefore,

\[
P=\mathrm{True}.
\]

The spam rule uses the OR operator:

\[
\boxed{P\lor Q}
\]

The message contains only two promotional keywords, so it does not contain
more than five.

Therefore,

\[
Q=\mathrm{False}.
\]

\end{exampleblock}

\end{frame}

% =========================================================
% ANSWER 2 -- PART 2
% =========================================================

\begin{frame}{Introduction -- Question 2: Answer (Continued)}

\begin{exampleblock}{Solution}

Substituting the truth values,

\[
P\lor Q
=
\mathrm{True}\lor\mathrm{False}
=
\mathrm{True}.
\]

Therefore, the message is classified as spam.

\medskip

Knowing only $Q$ is insufficient because the result of

\[
P\lor Q
\]

depends on both propositions.

\medskip

\textbf{Key idea:}

A disjunction is true when at least one of its propositions is true.

\end{exampleblock}

\end{frame}

% =========================================================
% QUESTION 3 -- PART 1
% =========================================================

\begin{frame}{Introduction -- Question 3}

\begin{block}{Database Access Control}

A database administrator defines the following access policy:

\begin{center}
\textit{A user can access a confidential file if the user is an administrator
or has been explicitly granted permission.}
\end{center}

Let

\[
A=\text{``the user is an administrator''}
\]

and

\[
P=\text{``the user has explicit permission.''}
\]

\end{block}

\end{frame}

% =========================================================
% QUESTION 3 -- PART 2
% =========================================================

\begin{frame}{Introduction -- Question 3 (Continued)}

\begin{block}{Question}

A user who is not an administrator and has not been granted permission
attempts to access the file.

\medskip

Determine whether access should be permitted.

\medskip

Represent the access policy using a logical expression.

\medskip

Explain how the truth values of $A$ and $P$ determine the result.

\end{block}

\end{frame}

% =========================================================
% ANSWER 3 -- PART 1
% =========================================================

\begin{frame}{Introduction -- Question 3: Answer}

\begin{exampleblock}{Solution}

The policy states that access is allowed when either condition is satisfied.

Therefore, the logical expression is

\[
\boxed{A\lor P}.
\]

The user is not an administrator, so

\[
A=\mathrm{False}.
\]

The user has not been granted permission, so

\[
P=\mathrm{False}.
\]

\end{exampleblock}

\end{frame}

% =========================================================
% ANSWER 3 -- PART 2
% =========================================================

\begin{frame}{Introduction -- Question 3: Answer (Continued)}

\begin{exampleblock}{Solution}

Substituting the truth values,

\[
A\lor P
=
\mathrm{False}\lor\mathrm{False}
=
\mathrm{False}.
\]

Therefore,

\[
\boxed{\text{Access is denied.}}
\]

\medskip

\textbf{Key idea:}

For a disjunction $A\lor P$ to be true, at least one of $A$ or $P$
must be true.

\end{exampleblock}

\end{frame}

% =========================================================
% QUESTION 4 -- PART 1
% =========================================================

\begin{frame}{Introduction -- Question 4}

\begin{block}{Machine Learning Model Evaluation}

A machine-learning model is designed to flag fraudulent transactions.

The system uses the statement:

\begin{center}
\textit{``If a transaction is flagged as suspicious, then it is fraudulent.''}
\end{center}

During testing, a transaction is flagged as suspicious, but is later found
to be legitimate.

\end{block}

\end{frame}

% =========================================================
% QUESTION 4 -- PART 2
% =========================================================

\begin{frame}{Introduction -- Question 4 (Continued)}

\begin{block}{Question}

Let

\[
P=\text{``the transaction is flagged as suspicious''}
\]

and

\[
Q=\text{``the transaction is fraudulent.''}
\]

\medskip

Determine whether the observation invalidates the stated implication.

Explain why or why not using the logical meaning of implication.

Give a condition under which the implication would definitely be false.

\end{block}

\end{frame}

% =========================================================
% ANSWER 4 -- PART 1
% =========================================================

\begin{frame}{Introduction -- Question 4: Answer}

\begin{exampleblock}{Solution}

The statement is represented as

\[
P\rightarrow Q.
\]

In the observed case, the transaction was flagged as suspicious.

Therefore,

\[
P=\mathrm{True}.
\]

However, the transaction was later found to be legitimate.

Therefore,

\[
Q=\mathrm{False}.
\]

\end{exampleblock}

\end{frame}

% =========================================================
% ANSWER 4 -- PART 2
% =========================================================

\begin{frame}{Introduction -- Question 4: Answer (Continued)}

\begin{exampleblock}{Solution}

An implication is false only in the case

\[
\mathrm{True}\rightarrow\mathrm{False}.
\]

Here,

\[
P\rightarrow Q
=
\mathrm{True}\rightarrow\mathrm{False}
=
\mathrm{False}.
\]

Therefore, the observation invalidates the stated implication.

\medskip

This is an example of a \textbf{false positive}: the system flags a
legitimate transaction as suspicious.

\end{exampleblock}

\end{frame}

% =========================================================
% QUESTION 5 -- PART 1
% =========================================================

\begin{frame}{Introduction -- Question 5}

\begin{block}{Autonomous System Decision}

An autonomous vehicle uses the following rule:

\begin{center}
\textit{If the road is wet and visibility is poor, the vehicle reduces
its speed.}
\end{center}

Let

\[
W=\text{``the road is wet''},
\]

\[
V=\text{``visibility is poor''},
\]

\[
S=\text{``the vehicle reduces speed.''}
\]

\end{block}

\end{frame}

% =========================================================
% QUESTION 5 -- PART 2
% =========================================================

\begin{frame}{Introduction -- Question 5 (Continued)}

\begin{block}{Question}

The following observations are recorded:

\begin{center}
\begin{tabular}{ll}
\toprule
\textbf{Condition} & \textbf{Observation}\\
\midrule
Road is wet & True\\
Visibility is poor & False\\
Vehicle reduces speed & False\\
\bottomrule
\end{tabular}
\end{center}

Represent the decision rule using $W$, $V$, and $S$.

Determine whether the observed situation violates the decision rule.

\end{block}

\end{frame}

% =========================================================
% ANSWER 5 -- PART 1
% =========================================================

\begin{frame}{Introduction -- Question 5: Answer}

\begin{exampleblock}{Solution}

The rule states:

\[
\boxed{(W\land V)\rightarrow S}
\]

From the observations,

\[
W=\mathrm{True},
\]

\[
V=\mathrm{False},
\]

\[
S=\mathrm{False}.
\]

Therefore,

\[
W\land V
=
\mathrm{True}\land\mathrm{False}
=
\mathrm{False}.
\]

\end{exampleblock}

\end{frame}

% =========================================================
% ANSWER 5 -- PART 2A
% =========================================================

\begin{frame}{Introduction -- Question 5: Answer (Continued)}

\begin{exampleblock}{Solution}

The complete implication becomes

\[
\mathrm{False}\rightarrow\mathrm{False}.
\]

An implication with a false antecedent is true.

Therefore, the observed situation does not violate the decision rule.

\medskip

A violating observation would be

\[
W=\mathrm{True},
\qquad
V=\mathrm{True},
\qquad
S=\mathrm{False}.
\]

\end{exampleblock}

\end{frame}

% =========================================================
% ANSWER 5 -- PART 2B
% =========================================================

\begin{frame}{Introduction -- Question 5: Answer (Final)}

\begin{exampleblock}{Solution}

For the violating observation,

\[
(W\land V)\rightarrow S
=
\mathrm{True}\rightarrow\mathrm{False}.
\]

Therefore,

\[
(W\land V)\rightarrow S=\mathrm{False}.
\]

\medskip

\textbf{Conclusion:}

The rule is violated when both required conditions occur, but the required
action does not occur.

\end{exampleblock}

\end{frame}

% =========================================================
% END OF TOPIC 1
% =========================================================
